\tableofcontents
\section{পরীক্ষা English}\label{sec:test}
\newcounter{custom}\setcounter{custom}{12}
Custom counter: \arabic{custom}; literal: 12 ১২।
\footnote{বাংলা footnote 12 ১২।}
\begin{enumerate}
\item প্রথম\label{item:test}
\begin{enumerate}\item দ্বিতীয় স্তর\item আরেকটি\end{enumerate}
\item দ্বিতীয়
\end{enumerate}
\begin{itemize}\item বাংলা bullet\end{itemize}
\begin{tabular}{lp{5cm}r}
\textbf{নাম}&\textbf{বিবরণ}&\textbf{সংখ্যা}\\
বাংলা&English এবং বাংলা&12\\
বাংলা&\textbf{গাঢ় লেখা}&১২
\end{tabular}
\begin{longtable}{ll}বাংলা&English\\অন্য সারি&১২\end{longtable}
\begin{figure}[ht]\centering\rule{2cm}{1cm}
\caption{বাংলা Figure}\label{fig:test}\end{figure}
\begin{equation}x=1\label{eq:structure}\end{equation}
\begin{subequations}\begin{align}a&=1\\b&=2\end{align}\end{subequations}
\begin{theorem}\label{thm:test}বাংলা theorem।\end{theorem}
\begin{proof}বাংলা proof।\end{proof}
\ref{sec:test}; \autoref{sec:test}; \autoref{fig:test};
\eqref{eq:structure}; \ref{item:test}; \autoref{thm:test}।
\appendix\section{পরিশিষ্ট পরীক্ষা}\label{app:test}
পরিশিষ্ট \ref{app:test}।
